%=============================================================================!
% 5.1  TURNING IN A PLANE                                                     !
%=============================================================================!
\section{Turning in a plane}
\label{sec:ro-plane}

A puck slides without friction on a smooth level table, an air table that
floats it on a thin cushion of air. A string tied to the puck runs to a
small hole in the middle of the table and through it to a hand underneath.
Set going sideways with the string taut, the puck circles the hole; pull
the string down, and the puck circles closer and faster. This section
finds the quantity that stays fixed while the speed changes, and the
reason it does.

\subsection*{Distance and angle}

Put the origin at the hole and describe the puck's position by its
distance $r$ from the hole and the angle $\theta$ that the line from the
hole makes with the $x$ axis, so that, by the definition of sine and
cosine scaled to a circle of radius $r$ (\cref{sec:mo-trig}),
\begin{equation*}
   x = r\cos\theta , \qquad y = r\sin\theta .
\end{equation*}
These are called polar coordinates. As the puck moves, both $r$ and
$\theta$ change with time, at the rates $\dot r$ and $\dot\theta$, the dot
marking the rate of change with time as in \cref{sec:mo-acceleration}.
Differentiating $x$ and $y$ with the product rule, and $\cos\theta$ and $\sin\theta$ with the chain rule, so that
$\dd\cos\theta/\dd t = -\sin\theta\,\dot\theta$ and $\dd\sin\theta/\dd t =
\cos\theta\,\dot\theta$,
\begin{equation*}
   \dot{x} = \dot{r}\cos\theta - r\dot\theta\sin\theta , \qquad
   \dot{y} = \dot{r}\sin\theta + r\dot\theta\cos\theta .
\end{equation*}
Two vectors of length 1, called unit vectors,
make this easier to read: $\vb{e}_r = (\cos\theta, \sin\theta)$, pointing
straight away from the hole, and $\vb{e}_\theta = (-\sin\theta,
\cos\theta)$, at right angles to it in the direction of increasing
$\theta$, anticlockwise; their scalar product is $-\cos\theta\sin\theta +
\sin\theta\cos\theta = 0$. Collecting the terms in $\dot r$ and in $r\dot
\theta$ in each component, $(\dot x, \dot y) = \dot r\,(\cos\theta,
\sin\theta) + r\dot\theta\,(-\sin\theta, \cos\theta)$, so the velocity is
\begin{equation}
   \vb{v} = \dot{r}\,\vb{e}_r + r\dot\theta\,\vb{e}_\theta ,
   \label{eq:ro-polar}
\end{equation}
a part $v_r = \dot r$ straight towards or away from the hole and a part
$v_\theta = r\dot\theta$ across the line to it. On a circle $r$ is
constant, $\dot\theta$ is the angular speed $\omega$, and $v_\theta =
r\dot\theta$ is the speed $R\omega$ of \cref{sec:mo-circle}.

\subsection*{Angular momentum and torque}

For a body of mass $m$ moving in the plane, define
\begin{equation}
   L = m\,(x v_y - y v_x) ,
   \label{eq:ro-ell}
\end{equation}
its angular momentum about the origin. In polar coordinates,
\begin{align*}
   x v_y - y v_x
   &= r\cos\theta\,(\dot{r}\sin\theta + r\dot\theta\cos\theta)
      && \text{putting in $x$, $y$, $\dot x$ and $\dot y$}\\
   &\quad - r\sin\theta\,(\dot{r}\cos\theta - r\dot\theta\sin\theta)\\
   &= r^2\dot\theta\,(\cos^2\theta + \sin^2\theta) = r^2\dot\theta
      && \text{the terms in $\dot r$ cancel; $\cos^2\theta + \sin^2\theta
      = 1$,}
\end{align*}
so $L = mr^2\dot\theta = mr\,v_\theta$: the distance from the origin times
the part of the momentum across the line to it. Moving straight towards
or away from the origin gives no angular momentum; going round it
anticlockwise gives a positive $L$, clockwise a negative one. Its unit is
\si{\kilogram\metre\squared\per\second}.

How does $L$ change? Differentiating \cref{eq:ro-ell} by the product rule,
\begin{align*}
   \frac{\dd L}{\dd t}
   &= m\,(\dot{x}v_y + x\dot{v}_y - \dot{y}v_x - y\dot{v}_x)\\
   &= m\,(x a_y - y a_x)
   && \text{since $\dot x v_y = v_xv_y = \dot y v_x$ and $\dot{\vb{v}} =
      \vb{a}$,}\\
   &= x F_y - y F_x
   && \text{by the second law, $m\vb{a} = \vb{F}$.}
\end{align*}
The combination on the right is the torque of the force about the origin,
\begin{equation}
   G = x F_y - y F_x , \qquad\text{and}\qquad
   \frac{\dd L}{\dd t} = G .
   \label{eq:ro-torque-plane}
\end{equation}
Torque is to angular momentum what force is to momentum. Writing the
force as $\vb{F} = F_r\vb{e}_r + F_\theta\vb{e}_\theta$, the same algebra
as for $L$, with $F_r$ in place of $\dot r$ and $F_\theta$ in place of
$r\dot\theta$, gives $G = rF_\theta$: only the part
of the force across the line from the origin turns the body, and it turns
it more the further out it acts. Its unit is the newton metre.

This is the lever. A door turns about its hinge (\cref{fig:ro-lever}a). A
push of \SI{10}{\newton} at right angles to the door at the handle,
\SI{0.8}{\metre} from the hinge, gives a torque of $0.8\times10 =
\SI{8}{\newton\metre}$; the same torque \SI{0.2}{\metre} from the hinge
needs $8/0.2 = \SI{40}{\newton}$, and a push straight towards the hinge,
with no part across the door, does not turn it at all.

Chapter~2 noted that two equal and opposite pushes at opposite corners of
a book set it spinning although they add to zero. Put the origin at the
centre of the book, with a push $(0, F)$ at the corner $(a, b)$ and a push
$(0, -F)$ at the opposite corner $(-a, -b)$ (\cref{fig:ro-lever}b). The
net force is zero, but the torques are $aF - b\times0 = aF$ and $(-a)(-F)
- (-b)\times0 = aF$, a total of $2aF$. Such a pair of forces, adding to no
force but to a torque, is called a couple.

\begin{figure}[htb]
\centering
\begin{tikzpicture}[line cap=round, line join=round,
   force/.style={-{Stealth[length=5pt]}, line width=1.0pt, accent}]
   % (a) a door seen from above, turning about its hinge
   \begin{scope}
      \fill[black!25] (-0.15,-0.35) rectangle (0,0.35);
      \draw[line width=2.2pt, black!60] (0,0) -- (3.2,0);
      \fill (0,0) circle (0.06);
      \node[left] at (-0.2,0) {\small hinge};
      \draw[force] (2.6,0) -- (2.6,1.1) node[above] {\small $\vb{F}$};
      \draw[{Stealth[length=3pt]}-{Stealth[length=3pt]}, line width=0.4pt]
         (0,-0.45) -- (2.6,-0.45) node[midway, below] {\small $r$};
      \draw[black!50, ->, line width=0.5pt] (1.4,0.25) arc[start angle=10,
         end angle=50, radius=1.4];
      \node at (0.15,1.6) {\small (a)};
   \end{scope}
   % (b) a book pushed at two opposite corners: a couple
   \begin{scope}[xshift=6.2cm, yshift=0.2cm]
      \draw[line width=0.6pt, fill=black!8] (-1.0,-0.7) rectangle (1.0,0.7);
      \fill (0,0) circle (0.04);
      \draw[force] (1.0,0.7) -- (1.0,1.5) node[above] {\small $(0, F)$};
      \draw[force] (-1.0,-0.7) -- (-1.0,-1.5) node[below] {\small $(0, -F)$};
      \node[right] at (1.0,0.55) {\small $(a, b)$};
      \node[left] at (-1.0,-0.55) {\small $(-a, -b)$};
      \node at (-1.9,1.4) {\small (b)};
   \end{scope}
\end{tikzpicture}
\caption{Torque. (a) A door seen from above: a force $\vb{F}$ at right
angles to the door, a distance $r$ from the hinge, has the torque $rF$
about the hinge. (b) A book pushed at two opposite corners: the forces add
to zero, but their torques about the centre add to $2aF$, and the book
spins.}
\label{fig:ro-lever}
\end{figure}

\subsection*{Pulling the string}

The string pulls the puck straight towards the hole, so its force has no
part across the line to the hole, $F_\theta = 0$, and its torque about the
hole is zero. The weight of the puck and the push of the table are
vertical and turn nothing in the plane. By \cref{eq:ro-torque-plane},
therefore, $L$ does not change: the angular momentum of the puck is
conserved, whatever the hand does with the string. A force that always
points towards or away from one fixed point is called a central force,
and every central force conserves the angular momentum about its centre.

Let the puck, of mass \SI{0.2}{\kilogram}, circle at $r =
\SI{0.5}{\metre}$ with speed \SI{1.2}{\metre\per\second}, so that $L = mr
v_\theta = 0.2\times0.5\times1.2 = \SI{0.12}{\kilogram\metre\squared\per
\second}$, and pull the string down at a steady \SI{0.05}{\metre\per
\second} until $r = \SI{0.25}{\metre}$. Since $L = mrv_\theta$ stays
fixed, $v_\theta = L/(mr)$ grows as $r$ shrinks: at \SI{0.25}{\metre} the
puck goes round at $0.12/(0.2\times0.25) = \SI{2.4}{\metre\per\second}$,
twice as fast at half the distance (\cref{fig:ro-puck}). Since
$\vb{e}_r$ and $\vb{e}_\theta$ are at right angles and of length 1,
Pythagoras' theorem applied to \cref{eq:ro-polar} gives $\lvert\vb{v}
\rvert^2 = \dot r^2 + v_\theta^2$, so the kinetic energy splits into a
part across the line, $\frac12 mv_\theta^2$, and a part along it, $\frac12
m\dot r^2$. The part across the line rises from $\frac12\times0.2\times
1.2^2 = \SI{0.144}{\joule}$ to $\frac12\times0.2\times2.4^2 =
\SI{0.576}{\joule}$, while the part along the line, with $\dot r =
\SI{-0.05}{\metre\per\second}$, stays a constant $\frac12\times0.2\times
0.05^2 = \SI{0.25}{\milli\joule}$. The extra
\SI{0.432}{\joule} is the work of the hand: by \cref{sec:en-kinetic} the
change in kinetic energy is the work done on the puck, the vertical forces
do none, and the string pulls inwards while the puck moves inwards, so its
work is positive.

\begin{figure}[htb]
\centering
\includegraphics{ch05_rotation/puck}
\caption{A puck of \SI{0.2}{\kilogram} on a string through a hole, circling
at \SI{0.5}{\metre} and \SI{1.2}{\metre\per\second} and then pulled in at
a steady \SI{0.05}{\metre\per\second} to \SI{0.25}{\metre}, computed by
stepping Newton's second law forward in short steps of time
(\texttt{solve\_newton} in \texttt{mdlab}). (a) The path seen from above, from the dot; the
shaded sectors are swept in \SI{0.25}{\second} each, at the start, the
middle and the end of the pull, and all three have the area
\SI{0.075}{\metre\squared}. (b) The angular momentum $L$, the speed and the
kinetic energy $K$, each divided by its value at the start, against the
distance from the hole.}
\label{fig:ro-puck}
\end{figure}

\subsection*{Equal areas in equal times}

The line from the hole to the puck sweeps out area as the puck moves. A
whole circle, of area $\pi r^2$, has the angle $2\pi$, so a sector of
angle $\Delta\theta$ has the area $\pi r^2\,\Delta\theta/(2\pi) = \frac12
r^2\,\Delta\theta$. In a short time $\Delta t$ the angle grows by about
$\dot\theta\,\Delta t$, and the line sweeps a thin sector of area $\Delta A
\approx \frac12 r^2\dot\theta\,\Delta t$; the change of $r$ in that time
alters the area only by a part proportional to $(\Delta t)^2$. Dividing by
$\Delta t$ and letting $\Delta t \to 0$, as in \cref{sec:mo-velocity},
\begin{equation*}
   \frac{\dd A}{\dd t} = \tfrac12 r^2\dot\theta = \frac{L}{2m}
   \qquad\text{since $L = mr^2\dot\theta$.}
\end{equation*}
While $L$ is constant the line sweeps equal areas in equal times, here
$0.12/(2\times0.2) = \SI{0.3}{\metre\squared}$ every second: long thin
sectors far out, short wide ones close in (\cref{fig:ro-puck}a). A planet
is pulled by the Sun's gravity towards the Sun, a central force, and the
same argument gives the law of equal areas that Kepler found in the
motions of the planets.

\begin{keyidea}
In a plane, the angular momentum $L = m(xv_y - yv_x) = mr v_\theta$
changes at the rate of the torque, $G = xF_y - yF_x = rF_\theta$. A
central force has no torque about its centre, so it conserves $L$: pulled
inwards, a body goes round faster.
\end{keyidea}

\notebook{05}{pull the puck in and out and watch $L$, the speed and the
energy}

\begin{selfcheck}
The puck circling at \SI{0.5}{\metre} and \SI{1.2}{\metre\per\second} is
let out slowly to \SI{1}{\metre}. How fast does it then go round, and is
the work done on it by the string positive or negative?
\end{selfcheck}
